\chapter{Introduction}
\label{sec:origins-intro}
% ===================================================================

The preceding parts established a progression: a physics derived
from the structure of a graph-structured lambda theory; an argument
that agents in massive populations are structurally prevented from
seeing the true ontology of their GSLT; and a proposal connecting
the hypercomputational tower above a GSLT to the distinction between
sentience and consciousness.

The present part turns to a question that is in some sense prior to
all three: \emph{how does life arise at all?}
We connect the framework to two of the most mathematically serious
contemporary proposals about the origins of life.
The first is \emph{Assembly Theory}, due to Cronin and
Walker~\cite{cronin2023assembly}, which assigns to each physical
object an assembly index and a copy number and proposes that
their joint magnitude is a biosignature.
The second is Eric Smith's proposal~\cite{smith2016origin} that
life constitutes a \emph{fourth geosphere}, emerging from prebiotic
chemistry through a phase transition driven by autocatalytic
self-reinforcement.

The connections we draw are not analogies.
Assembly index is identified precisely with causal depth in the
synchronization tree of a GSLT.
Copy number is given a rigorous definition in terms of the
Rho calculus namespace structure, filling a gap in Assembly
Theory that has attracted criticism.
The phase transition is identified with the crossing of a basin
of attraction boundary in the logical topology of Part~II,
with the attractor being the concurrent fixed point --- the
Rho calculus Y combinator --- that is the abstract core of
hereditary replication.

As throughout this volume, the arguments are presented as
mathematical skeletons with gaps identified explicitly.

\section{Organization}

Section~\ref{sec:assembly} develops the connection to Assembly
Theory: assembly index as causal depth, copy number in a namespace,
and the biosignature criterion as a well-posed mathematical
statement.
Section~\ref{sec:replication} identifies replication as a fixed
point of the interaction dynamics and notes the proximity of the
comm rule to a genetic crossover operator.
Section~\ref{sec:phase} develops the connection to Smith's phase
transition via the density of recursive attractors in the logical
topology.
Section~\ref{sec:origins-gaps} catalogs the open gaps.

% ===================================================================
